\chapter{Discussion}
\label{chap:discussion}

\section{Summary of Findings}

The work pursued a single operational question: can a heterogeneous
knowledge graph, built from Indian court judgments and filtered to be
structurally leakage-safe, support an outcome-prediction model whose
accuracy is attributable to legal reasoning rather than to identity or
decision shortcuts? The chapters above produced four findings that,
taken together, answer that question affirmatively but with
important qualifications.

\paragraph{The pipeline produces a legally coherent corpus.} Across all
five legal-domain buckets, the top-PageRank entities recovered by
OpenNyAI NER and the canonicalisation stage are exactly the statutes,
provisions, and courts a lawyer would expect: IPC and CrPC in criminal
matters, the Motor Vehicles Act~1988 and Section~166 in motor-accident
matters, the Land Acquisition Act~1894 and the Constitution of India
in land-property matters, and the POCSO Act in sexual-offences matters
(Chapter~\ref{chap:graph_exploration}). The corpus is
\emph{legally faithful}, not merely statistically clean.

{Equally important, the reasoning-focused pruning introduced
in Chapter~\ref{chap:kg} does not collapse the graph into something too
thin to learn from. Removing \textsc{GPE}/\textsc{Org}/\textsc{Date}
nodes and forbidding shortcut edges from local identity proxies to the
case representation still leaves a large, connected, and legally
meaningful structure. That is a substantive design result in its own
right: the leakage controls are restrictive without being destructive.\par}

\paragraph{The graph is hub-dominated and cross-domain.} The $880$k-node
cross-bucket graph is organised around a small set of super-hubs --
IPC, CrPC, the Supreme Court, the Apex Court -- each of which bridges
all five legal domains. The four IPC/CrPC-driven buckets share a
significant fraction of entities with each other; land-property
partially detaches via its own statutory backbone. The cross-domain
transfer setting this thesis studies is, by construction, the right
setting for this graph.

\paragraph{The HGT model learns cross-case structural signal.} A
probing classifier trained on the $64$-dimensional post-GNN case
embedding reaches $0.7840$ accuracy on the cross-bucket test set
(Chapter~\ref{chap:explainability},
Table~\ref{tab:probing_results}), compared to $0.7443$ for the
pre-GNN raw features. The $+3.97$-point gain at fixed probing
capacity, with a $16\times$ smaller representation, isolates the
HGT's contribution from the underlying \texttt{bge-m3} text
features. The best end-to-end configuration
(\texttt{party\_args\_lr\_decay}) reaches $0.8020$ accuracy,
$0.7959$ macro-F1, and $0.8831$ ROC-AUC on the pooled five-domain
test set.

{The updated ablations also clarify \emph{what kind} of
gain this is. The graph does help, but not primarily by replacing text
with a pure hub-traversal signal: \texttt{text\_only} remains strong,
\texttt{no\_cross\_case} is close to baseline on the pooled task, and
the t-SNE and probing analyses show that message passing mainly
reorganises the already-strong text representation into a more
outcome-informative case embedding. In this regime, the graph's
contribution is best described as structured aggregation of legal
context, with cross-case shared authorities providing a secondary,
domain-dependent benefit rather than the entire predictive signal.\par}

\paragraph{The model does not rely on identity shortcuts.} The
\texttt{no\_names} ablation -- which strips party, lawyer, and judge
name content -- matches or slightly improves over the baseline on
every bucket. The t-SNE geometry of the \texttt{no\_names} embedding
is qualitatively indistinguishable from the baseline's, and its
probing accuracy is within $0.2$ points of the baseline probe. The
leakage-prevention strategy of
Section~\ref{sec:leakage_prevention} -- field-level removal,
rhetorical-role filtering, phrase-level masking, temporal gating on
shared entities, and keeping identity nodes local -- is both
necessary (without RR filtering the model could read the decision
verbatim) and sufficient (with all gates in place, removing identity
features is a no-op).

{The updated error analysis sharpens that conclusion. The
remaining difficult cases are not those where names have been removed;
they are the cases with the densest statute neighbourhoods and the
longest factual narratives. That same pattern appears in the baseline
and \texttt{no\_names} variants. The residual errors therefore look
more like failures on doctrinally crowded or document-heavy judgments
than failures caused by withholding identity information.\par}

{Taken together, the evidence for resisting identity-based
leakage shortcuts is now multi-layered rather than resting on a single
ablation. At the schema level, courts, judges, and lawyers are kept
local and shortcut edges from identity proxies are forbidden; at the
text level, rhetorical-role and phrase filtering remove direct decision
leakage; at the model level, \texttt{no\_names} matches baseline; and
at the representation level, probing, t-SNE, saliency, and error
analysis all point to the same conclusion: the learned signal is driven
primarily by legal text structure and shared authority nodes, not by
recoverable person- or bench-specific identity cues.\par}

\paragraph{Predictions are auditable at the case level.}
Chapter~\ref{chap:explainability} builds a four-level
interpretability pipeline over the frozen HGT: representation-level
probing and t-SNE, component-level node-zeroing and attention, error
decomposition, and the per-case
\path{Graph_Analyser} pipeline that combines a heterogeneous
PGExplainer, FAISS retrieval in post-GNN space, and a Mistral-Small
verbaliser that produces a legally named narrative grounded in the
extracted evidence bundle.  On the worked correct example the
pipeline correctly surfaces \emph{Section~320 CrPC} and
\emph{Section~482 CrPC} as the governing provisions and retrieves
three structurally near-identical Punjab quashing cases as analogues
--- exactly the legal scaffolding one would expect a reviewer to
check.

{More importantly, the same retrieval mechanism exposes a
sharp structural signature of the model's confident errors.  Of
1{,}532 high-confidence predictions, 72 are confidently wrong; for
those 72 cases the top-$5$ FAISS neighbours match the \emph{predicted}
(incorrect) label with mean frequency $0.989$, and the \emph{true}
label with mean frequency $0.011$.  A confident error on this model
is therefore not a random slip --- it is a case placed inside the
wrong outcome neighbourhood, whose prediction \emph{and} whose
retrieval-based explanation reinforce each other.  This is an
operationally important finding, because it means the post-hoc
explanations are honest in the sense that matters most: they are
willing to confirm a mistake rather than mask it, and a
class-boundary proximity score --- not classifier confidence ---
is the right deployment-time filter.\par}

\paragraph{The cross-bucket representation transfers, partially, to
an unseen legal domain.} Section~\ref{sec:cross_domain_food_safety}
evaluates the trained cross-bucket checkpoints zero-shot on
$11{,}769$ food-safety judgments --- a body of case law governed by
statutes none of the five training buckets contain.  Without any
fine-tuning, the model reaches $0.608\pm0.005$ accuracy,
$0.603\pm0.007$ macro-F1, and $0.675\pm0.020$ ROC-AUC against a
$0.543$ majority-class baseline.  The $20$-point gap to in-domain
performance is substantial and makes the regime's difficulty
explicit, but the strict separation from baseline, together with the
low across-fold variance, is evidence that the cross-bucket HGT
learns a \emph{partially} domain-general legal-reasoning prior
(argument-to-outcome mapping, petitioner/respondent dynamics,
procedural markers) rather than merely memorising the five training
domains' statutory hubs.  In combination with the leakage-safety
story, this is the clearest single indication that the gain
reported by the in-domain tables is not an artefact of in-sample
corpus regularities.

\section{Limitations}

Several limitations bound the interpretation of these results.

\paragraph{Transductive evaluation assumes entity overlap at test
time.} The main Chapter~\ref{chap:results} numbers are produced
under the transductive protocol, i.e.\ the full graph (including
test cases' entity neighbourhoods) is visible during training.  This
matches the realistic deployment setting (a new judgment is appended
to an existing graph built from the same canonical vocabulary).  The
food-safety experiment in
Section~\ref{sec:cross_domain_food_safety} partially relaxes this by
evaluating on an \emph{entirely unseen legal domain} with disjoint
statutes and provisions, and the $+6.5$-point margin over the
majority baseline there shows that at least some of the learned
signal survives a hard entity shift.  What that experiment does
\emph{not} measure is the finer-grained inductive generalisation
within a familiar domain --- e.g.\ a calendar-year or court
hold-out, in which most entities are known but the specific case
neighbourhood is new.  The per-year error analysis already suggests
older cases are harder, which is a prior on how such a split would
behave, and it remains the natural follow-on experiment.

\paragraph{Outcome labels are LLM-generated.} The win/loss labels are
produced by prompting a Mistral model on the OpenNyAI-produced summary
of each case. The labels are therefore noisy with respect to the
ground-truth disposition, and the noise is not uniform across
buckets: sexual-offences and financial-fraud matters often produce
multi-part dispositions that do not map cleanly onto a binary
win/loss axis. The fact that these are also the two worst-performing
buckets is at least partially a label-noise effect rather than a
modelling failure. A multi-label or ordinal target is produced by the
same pipeline and is an obvious next experiment.

\paragraph{RPC/RLC filtering is a lower bound on leakage, not a
guarantee.} The leakage-safety story combines a deterministic filter
(rhetorical-role labels, field names, leakage phrases) with a
structural argument (shared-entity types are law, not identity).
Neither is complete: an RR classifier can mis-label a decisive
sentence, a leakage phrase list is never exhaustive, and a sufficiently
expressive GNN could in principle reconstruct a decisive fragment
from non-decisive neighbours. The \texttt{no\_names} result is
evidence against one specific leakage mode (party identity) but not
against all leakage modes. A manual audit of high-confidence correct
predictions -- sampled across buckets and rhetorical-role density --
is the next concrete step.

\paragraph{The ablation ordering is stable but the margins are
small.} The best configuration beats the baseline by $+1.7$ points of
accuracy on the pooled dataset; most other variants are within
$\pm 1$ point. The 5-fold error bars are of the same order, which
means that while the \emph{qualitative} findings (text dominates,
identity is removable, motor-accidents is easiest) are robust, the
\emph{quantitative} ordering of the close ablations should be
reported as a cluster of comparable configurations rather than a
strict ranking.

{\paragraph{The pooled optimum is not the universal optimum.}
The new per-dataset best-variant table in Chapter~\ref{chap:results}
shows that the strongest global configuration is not the best choice in
every domain: family-matrimonial, motor-accidents, and
sexual-offences each favour different ablations under accuracy or
macro-F1. This means the pooled cross-bucket score should be read as a
useful compromise over heterogeneous graph geometries, not as evidence
that one architectural setting is intrinsically best for all legal
subdomains.}

\section{Future Work}

Several directions follow naturally from the findings and limitations
above, in rough order of expected payoff.

\paragraph{1.\ Inductive and temporal splits.} The transductive
protocol is the easier regime; a calendar-year hold-out (train on
pre-2020, test on post-2020) and a court hold-out (train on all but
one High Court, test on that one) would measure inductive
generalisation directly. The connectivity-score machinery in
Chapter~\ref{chap:graph_exploration} already provides the right
diagnostic for these splits: a well-generalising model should remain
competitive on the hold-out distribution's low-connectivity tail.

\paragraph{2.\ Multi-label and ordinal outcome targets.} The Mistral
labelling pipeline produces a multi-label variant alongside the
single-label \texttt{win}/\texttt{loss} target used in this thesis.
Moving the classifier head to a multi-label sigmoid (or to an ordinal
regression over the $[-1,+1]$ score) would (i) absorb the
disposition ambiguity currently pushed into the binary target and
(ii) enable disposition-specific error analysis (e.g.\ \emph{remanded}
cases separately from \emph{dismissed} cases), which the present
evaluation cannot distinguish.

\paragraph{3.\ Expert-in-the-loop validation of the case-level
explanations.} Chapter~\ref{chap:explainability} already produces
per-case explanations: the \path{Graph_Analyser} pipeline extracts,
for any test case, the top statutes, provisions, precedents, and
argument spans the frozen HGT relied on, together with the
FAISS-retrieved similar training cases and an LLM-written narrative
grounded in that evidence bundle.  The natural next step is the
validation loop this machinery enables: sample predictions across
confidence and difficulty strata, hand the Phase~4 evidence bundle
and the Phase~5 narrative to a legal expert, and collect structured
judgements on (i) whether the surfaced statutes and provisions are
legally defensible reasons for the predicted outcome, (ii) whether
the retrieved analogues are \emph{legitimate} analogues rather than
surface lookalikes, and (iii) whether the confident-wrong neighbour
pattern identified in Chapter~\ref{chap:explainability} corresponds
to a coherent class of distinguishing-fact omissions.  That turns
the explainer from an internal diagnostic into an auditable
artefact and is the cleanest path from the present offline
evaluation to a deployable review tool.

\paragraph{4.\ Scaling the bucket inventory.} Five legal domains is
a deliberate scope; the pipeline is bucket-agnostic by construction,
and adding new domains (tax, labour, consumer) would expand the
cross-domain spine and, through shared IPC/CrPC hubs, improve signal
on the existing buckets as well. The main engineering cost is the
per-bucket Mistral labelling pass; the rest of the pipeline scales
without modification.

\paragraph{5.\ Stronger identity-leakage audits.} The \texttt{no\_names}
ablation is a single, coarse leakage probe. Finer probes -- a
judge-hold-out split to measure judge-identity leakage directly; a
court-hold-out split to measure court-identity leakage; and a
text-similarity-based nearest-neighbour attack to detect cases whose
neighbours in the graph also share rare token n-grams -- would
tighten the leakage-safety claim from ``no evidence of identity
shortcuts'' to ``identity shortcuts are quantitatively bounded''.

\paragraph{6.\ Cross-domain fine-tuning and domain-adaptive
heads.}  The food-safety experiment in
Section~\ref{sec:cross_domain_food_safety} establishes the
cross-bucket HGT as a usable zero-shot base model on unseen Indian
legal domains, but also quantifies the $\sim\!20$-point in-domain
vs.\ out-of-domain gap.  Two concrete follow-ons map directly onto
that finding: (i) light-weight fine-tuning --- freezing the HGT
trunk and re-training only the classifier head (and optionally a
small adapter over the \textsc{Case} node) on a modestly-sized
target-domain labelled sample, to quantify how quickly the gap
closes as a function of labelled food-safety examples; and
(ii) a domain-conditioned head that reads a domain embedding
alongside the post-GNN case embedding, so that one unified model
can be evaluated across all training domains \emph{and} a held-out
new domain without per-domain retraining.  Both are natural
complements to the zero-shot baseline reported here.

{\paragraph{7.\ Connectivity- and difficulty-aware modelling.}
Two empirical regularities now point to a targeted next step. First,
the pooled best ablation is not the best ablation inside every bucket.
Second, the hardest held-out cases are the ones with long factual
narratives or dense statute neighbourhoods. A natural extension is
therefore a model that allocates capacity conditionally -- for example,
using richer hierarchical text encoding for document-heavy cases,
stronger graph propagation for high-connectivity cases, or a
domain-adaptive head once the bucket geometry is known. That would turn
the descriptive error analysis of Chapter~\ref{chap:results} into an
explicit modelling strategy.}

\medskip

Taken together, these findings and directions support a narrow and
defensible thesis: legal outcome prediction over Indian court
judgments \emph{is} improvable by a heterogeneous, cross-case graph
representation; the improvement is modest but measurable; it survives
an aggressive identity-leakage audit; it produces predictions that
can be audited case-by-case with legally named evidence; and it
transfers partially, without fine-tuning, to a legal domain the
model has never seen.

{More strongly, the combination of schema design,
leakage-filtered text construction, locality constraints on identity
nodes, post-hoc ablation evidence, per-case interpretability via the
\path{Graph_Analyser} pipeline, and a zero-shot out-of-domain test
on food-safety judgments shows that this framework does not merely
claim to avoid identity shortcuts and produce a usable prediction
number; it operationally resists leakage while remaining predictive,
it exposes its own confident failure mode through representation-space
neighbour analysis, and it carries a non-trivial legal-reasoning
prior across domain boundaries. That combination --- leakage-safe
construction, auditable predictions, and measured cross-domain
transfer --- is the main methodological contribution of the thesis,
not just a side condition on the experimental setup.}
